\documentclass[10pt,a4paper]{amsart}
\usepackage[margin=19mm]{geometry}
\usepackage{amsmath,amssymb,mathtools}
\usepackage[scr=rsfs,cal=euler]{mathalfa}
\usepackage{xfrac}
\usepackage[unicode,hidelinks,bookmarks=false]{hyperref}
\newcommand{\ie}{i.e.\ }
\newcommand{\cf}{cf.\ }
\newcommand{\resp}{respectively\ }
\newcommand{\Subsection}{\subsection}
\providecommand{\nobreakdash}{\nobreak\hbox{-}}
\setlength{\emergencystretch}{2em}
\multlinegap0pt
\usepackage{float}
\usepackage{amssymb}
\usepackage{algorithm}\usepackage{algpseudocode}
\usepackage[shortlabels]{enumitem}
\usepackage{tikz}
\newtheorem{theorem}{Theorem}
\title{Data dependent Shepard approximation through and adaptive modification of the shape parameter}
\author{Jos\'e Kuruc, Juan Ruiz-\'Alvarez, Bo Wang, Dionisio-F\'elix Y\'a\~nez}
\date{Source: arXiv:2606.20332v1 (CC BY 4.0)}
\begin{document}
\maketitle

\noindent\emph{Source and licence.} Data dependent Shepard approximation through and adaptive modification of the shape parameter. Version arXiv:2606.20332v1 (CC BY 4.0), distributed by arXiv under the Creative Commons Attribution 4.0 International licence, http://creativecommons.org/licenses/by/4.0/, as recorded in the arXiv licence metadata for this exact version and shown on the abstract page https://arxiv.org/abs/2606.20332v1. Full licence notice: \texttt{licenses/arxiv-2606.20332-CC-BY-4.0.txt}.\par\smallskip

\noindent\textit{Editorial scope.} Counted unit: the article's theorem thm:thinning-local (source0.tex lines 360--398). The article's complete original proof of it is delivered with it (source0.tex lines 400--471, one \texttt{proof} environment of 72 lines). No mathematics is written, repaired or re-derived: the statement and the proof are the article's own lines, in the article's own notation, and no retained line is substituted or re-flowed.  Retained uncounted as the source-local context the proof needs: lines 275--279; lines 347--350; lines 354--357.  Nothing was omitted from any retained span, so the package carries no omission record.  No retained line contains a citation, so the wrapper carries no bibliography and no bibliography entry.  No retained line contains a figure environment, an image-inclusion command or drawing code, so no image is delivered and the package declares no asset dependency.  Editorial formatting only, in addition: an AMS \texttt{amsart} preamble that restates the article's own declarations and macros used by the retained lines, namely the environments \texttt{theorem} on separate counters, together with the standard packages those lines need.  The retained environments print in this document as Theorem 1 in order of appearance, whereas the article prints the same items with the numbering of its own full document.  The complete native source of the article as distributed in the arXiv e-print for this exact version is bundled unchanged as \texttt{sources/203208/source0.tex}.
\begin{equation}\label{eq:error-split}
E(\mathbf{x}) := s(\mathbf{x})-f^{+}(\mathbf{x})
= \underbrace{\omega^{+}(\mathbf{x})(\bar f^{+}(\mathbf{x})-f^{+}(\mathbf{x}))}_{\text{smooth-side interpolation}}
+ \underbrace{\omega^{-}(\mathbf{x})(\bar f^{-}(\mathbf{x})-f^{+}(\mathbf{x}))}_{\text{wrong-side leakage}}.\
\end{equation}

\begin{equation}\label{salto}
\|f\|_{d}
:= \max\{\bigl| f^+(\mathbf{x}_i) - f^-(\mathbf{x}_j) \bigr|: i\in I^+_{\mathrm{near}},\,\,j \in I^-_{\mathrm{near}}\}.
\end{equation}

\begin{equation}\label{eq:Rdelta-local}
\delta\ >\ \delta_{\mathrm{far}}(\mathbf{x})+\delta_{\mathrm{smooth}}(\mathbf{x})
\quad\text{for all $\mathbf{x}$ in the near belt}.
\end{equation}

\begin{theorem}[The smearing belt around the discontinuity can be reduced]
\label{thm:thinning-local}
Under \textnormal{(A1)}-\textnormal{(A3)} and a tolerance \eqref{eq:Rdelta-local}, we define the
distance thresholds
\[
d_{\star}^{\mathrm{const}}(\mathbf{x})
:=\frac{1}{\varepsilon}\left[
\frac{1}{C_2}\log\!\Big(
\frac{A\,\|f\|_{d}}{
\delta\,\phi(\varepsilon\,h_{X^{+}}(\mathbf{x}))-\|f\|_{d}\,B_{\mathrm{far}}}
\Big)\right]^{\!1/p},
\]
\[
d_{\star}^{\mathrm{adap}}(\mathbf{x})
:=\frac{1}{\kappa\,\varepsilon}\left[
\frac{1}{C_2}\log\!\Big(
\frac{A\,\|f\|_{d}}{
\delta\,\phi(\varepsilon_{\mathrm{den}}(\mathbf{x})\,h_{X^{+}}(\mathbf{x}))
-\|f\|_{d}\,B_{\mathrm{far}}}
\Big)\right]^{\!1/p}.
\]
Then for each method (constant/adaptive), \emph{if} $d(\mathbf{x})\ge d_{\star}^{\mathrm{method}}(\mathbf{x})$
we have $|E(\mathbf{x})|\le \delta$. Consequently, the $\delta$-smearing belt in $\Omega^{+}$ is contained in
a tube of half-width $\sup_{\mathbf{x}\text{ on the belt}} d_{\star}^{\mathrm{method}}(\mathbf{x})$.
Moreover,
\[
\frac{\operatorname{width}(B_{\delta}^{\mathrm{adap}})}
     {\operatorname{width}(B_{\delta}^{\mathrm{const}})}
\ \le\
\frac{1}{\kappa}\left(
\frac{
\log\Big(\!\dfrac{A\,\|f\|_{d}}{\delta\,\phi(\varepsilon_{\mathrm{den}}\,h_{X^{+}}(\mathbf{x}))-\|f\|_{d}\,B_{\mathrm{far}}}\Big)
}{
\log\Big(\!\dfrac{A\,\|f\|_{d}}{\delta\,\phi(\varepsilon\,h_{X^{+}}(\mathbf{x}))-\|f\|_{d}\,B_{\mathrm{far}}}\Big)
}
\right)^{\!1/p},
\]
so the adaptive belt can be thinner by about a factor $1/\kappa$ (the logarithmic ratio is close to one).
\end{theorem}

\begin{proof}
We start from the exact identity and isolate the leakage.
From \eqref{eq:error-split} and $|\omega^{+}(\mathbf{x})(\bar f^{+}(\mathbf{x})-f^{+}(\mathbf{x}))|\le \delta_{\mathrm{smooth}}(\mathbf{x})$,
\[
|E(\mathbf{x})|\ \le\ \delta_{\mathrm{smooth}}(\mathbf{x})
+ \omega^{-}(\mathbf{x})\,|\bar f^{-}(\mathbf{x})-f^{+}(\mathbf{x})|
\ \le\ \delta_{\mathrm{smooth}}(\mathbf{x})
+ \frac{W^{-}(\mathbf{x})}{W^{+}(\mathbf{x})}\,\|f\|_{d},
\]
as by definition of $\|f\|_{d}$ in (\ref{salto}),
$|\bar f^{-}(\mathbf{x})-f^{+}(\mathbf{x})|\le \|f\|_{d}$ in the near belt.

\smallskip
Now, we insert the denominator lower bound and split $W^{-}$.
Using (B3), $W^{+}(\mathbf{x})\ge \phi_{\mathrm{den}}(\mathbf{x})
=\phi(\varepsilon_{\mathrm{den}}\,h_{X^{+}}(\mathbf{x}))$,
and the near/far split $W^{-}(\mathbf{x})=W^{-}_{\mathrm{near}}(\mathbf{x})+W^{-}_{\mathrm{far}}(\mathbf{x})$,
\[
|E(\mathbf{x})|\ \le\ \delta_{\mathrm{smooth}}(\mathbf{x})
+ \frac{\|f\|_{d}}{\phi_{\mathrm{den}}(\mathbf{x})}
\Big(W^{-}_{\mathrm{near}}(\mathbf{x})+W^{-}_{\mathrm{far}}(\mathbf{x})\Big).
\]
\smallskip
By \eqref{eq:Rdelta-local}, the \emph{gap}
\[
\Gamma(\mathbf{x})
:=\delta\,\phi_{\mathrm{den}}(\mathbf{x})
-\|f\|_{d}\,B_{\mathrm{far}}
-\phi_{\mathrm{den}}(\mathbf{x})\,\delta_{\mathrm{smooth}}(\mathbf{x})
\]
is strictly positive (the tolerance exceeds the floors). Therefore a \emph{sufficient} condition for
$|E(\mathbf{x})|\le \delta$ is
\[
  \delta_{\mathrm{smooth}}(\mathbf{x})
+ \frac{\|f\|_{d}}{\phi_{\mathrm{den}}(\mathbf{x})}
\Big(W^{-}_{\mathrm{near}}(\mathbf{x})+W^{-}_{\mathrm{far}}(\mathbf{x})\Big)\le\delta,
\]
and
\[
W^{-}_{\mathrm{near}}(\mathbf{x})\ \le\ \frac{\Gamma(\mathbf{x})}{\|f\|_{d}}
\ \le\ \frac{\delta\,\phi_{\mathrm{den}}(\mathbf{x})-\|f\|_{d}\,B_{\mathrm{far}}}{\|f\|_{d}}.
\]
where we have dropped $\phi_{\mathrm{den}}(\mathbf{x})\delta_{\mathrm{smooth}}(\mathbf{x})$, that is a positive quantity.

\smallskip
Now we use the near bound and solve for the distance $d$. By (B1), in the constant-$\varepsilon$ case,
$W^{-}_{\mathrm{near}}(\mathbf{x})\le A\,e^{-C_2(\varepsilon d)^p}$, hence it suffices that
\[
A\,e^{-C_2(\varepsilon d)^p}
\ \le\
\frac{\delta\,\phi(\varepsilon\,h_{X^{+}}(\mathbf{x}))-\|f\|_{d}\,B_{\mathrm{far}}}{\|f\|_{d}}.
\]
Taking logs and solving for $d$ gives
\[
d\ \ge\
\frac{1}{\varepsilon}\left[
\frac{1}{C_2}\log\!\Big(
\frac{A\,\|f\|_{d}}{
\delta\,\phi(\varepsilon\,h_{X^{+}}(\mathbf{x}))-\|f\|_{d}\,B_{\mathrm{far}}}
\Big)\right]^{\!1/p}
=:d_{\star}^{\mathrm{const}}(\mathbf{x}).
\]
The adaptive case is identical, replacing $\varepsilon$ by $\kappa\varepsilon$ (by (A3)) and
$\phi(\varepsilon h_{X^{+}})$ by $\phi(\varepsilon_{\mathrm{den}}(\mathbf{x}) h_{X^{+}})$, yielding
$d_{\star}^{\mathrm{adap}}(\mathbf{x})$.

\smallskip

Now, if $d(\mathbf{x})\ge d_{\star}^{\mathrm{method}}(\mathbf{x})$ then $|E(\mathbf{x})|\le \delta$, so the belt
is contained in the tube $\{d(\mathbf{x})<\sup d_{\star}^{\mathrm{method}}(\mathbf{x})\}$. This gives the
stated width bound and the ratio inequality, with the explicit $1/\kappa$ factor.
\end{proof}

\end{document}
